\section{Support and dimension of modules}
\label{section-support}

\noindent
Some basic results on the support and dimension of modules.

\begin{lemma}
\label{lemma-filter-Noetherian-module}
Let $R$ be a Noetherian ring, and let $M$ be a finite $R$-module.
There exists a filtration by $R$-submodules
$$
0 = M_0 \subset M_1 \subset \ldots \subset M_n = M
$$
such that each quotient $M_i/M_{i-1}$ is isomorphic
to $R/\mathfrak p_i$ for some prime ideal $\mathfrak p_i$
of $R$.
\end{lemma}

\begin{proof}[First proof]
By Lemma \ref{lemma-trivial-filter-finite-module}
it suffices to do the case $M = R/I$ for some ideal $I$.
Consider the set $S$ of ideals $J$ such that the lemma
does not hold for the module $R/J$, and order it by
inclusion. To arrive at a
contradiction, assume that $S$ is not empty. Because
$R$ is Noetherian, $S$ has a maximal element $J$.
By definition of $S$, the ideal $J$ cannot be prime.
Pick $a, b\in R$ such that $ab \in J$, but neither
$a \in J$ nor $b\in J$. Consider the filtration
$0 \subset aR/(J \cap aR) \subset R/J$.
Note that both the submodule $aR/(J \cap aR)$ and the quotient
module $(R/J)/(aR/(J \cap aR))$ are cyclic modules; write
them as $R/J'$ and $R/J''$ so we have a short exact sequence
$0 \to R/J' \to R/J \to R/J'' \to 0$.
The inclusion $J \subset J'$ is strict
as $b \in J'$ and the inclusion $J \subset J''$ is strict as $a \in J''$.
Hence by maximality of $J$, both $R/J'$ and $R/J''$ have a filtration as
above and hence so does $R/J$. Contradiction.
\end{proof}

\begin{proof}[Second proof]
For an $R$-module $M$ we say $P(M)$ holds if there exists a filtration
as in the statement of the lemma. Observe that $P$ is stable under
extensions and holds for $0$. By Lemma \ref{lemma-trivial-filter-finite-module}
it suffices to prove $P(R/I)$ holds for every ideal $I$.
If not then because $R$ is Noetherian, there is a maximal
counter example $J$. By Example \ref{example-oka-family-property-modules} and
Proposition \ref{proposition-oka}
the ideal $J$ is prime which is a contradiction.
\end{proof}

\begin{lemma}
\label{lemma-filter-primes-in-support}
Let $R$, $M$, $M_i$, $\mathfrak p_i$ as in
Lemma \ref{lemma-filter-Noetherian-module}.
Then $\text{Supp}(M) = \bigcup V(\mathfrak p_i)$
and in particular $\mathfrak p_i \in \text{Supp}(M)$.
\end{lemma}

\begin{proof}
This follows from Lemmas \ref{lemma-support-closed} and
\ref{lemma-support-quotient}.
\end{proof}

\begin{lemma}
\label{lemma-support-point}
Suppose that $R$ is a Noetherian local ring with
maximal ideal $\mathfrak m$. Let $M$ be a nonzero finite
$R$-module. Then $\text{Supp}(M) = \{ \mathfrak m\}$
if and only if $M$ has finite length over $R$.
\end{lemma}

\begin{proof}
Assume that $\text{Supp}(M) = \{ \mathfrak m\}$.
It suffices to show that all the primes $\mathfrak p_i$
in the filtration of Lemma \ref{lemma-filter-Noetherian-module}
are the maximal ideal. This is clear by
Lemma \ref{lemma-filter-primes-in-support}.

\medskip\noindent
Suppose that $M$ has finite length over $R$.
Then $\mathfrak m^n M = 0$ by Lemma \ref{lemma-length-infinite}.
Since some element of $\mathfrak m$ maps to a unit
in $R_{\mathfrak p}$ for any prime
$\mathfrak p \not = \mathfrak m$ in $R$ we see $M_{\mathfrak p} = 0$.
\end{proof}

\begin{lemma}
\label{lemma-Noetherian-power-ideal-kills-module}
Let $R$ be a Noetherian ring.
Let $I \subset R$ be an ideal.
Let $M$ be a finite $R$-module.
Then $I^nM = 0$ for some $n \geq 0$ if and only if
$\text{Supp}(M) \subset V(I)$.
\end{lemma}

\begin{proof}
Indeed, $I^nM = 0$ is equivalent to $I^n \subset \text{Ann}(M)$.
Since $R$ is Noetherian, this is equivalent to
$I \subset \sqrt{\text{Ann}(M)}$, see
Lemma \ref{lemma-Noetherian-power}.
This in turn is equivalent to $V(I) \supset V(\text{Ann}(M))$, see
Lemma \ref{lemma-Zariski-topology}.
By Lemma \ref{lemma-support-closed}
this is equivalent to $V(I) \supset \text{Supp}(M)$.
\end{proof}

\begin{lemma}
\label{lemma-filter-minimal-primes-in-support}
Let $R$, $M$, $M_i$, $\mathfrak p_i$ as in
Lemma \ref{lemma-filter-Noetherian-module}.
The minimal elements of the set $\{\mathfrak p_i\}$
are the minimal elements of $\text{Supp}(M)$.
The number of times a minimal prime $\mathfrak p$
occurs is
$$
\#\{i \mid \mathfrak p_i = \mathfrak p\}
=
\text{length}_{R_\mathfrak p} M_{\mathfrak p}.
$$
\end{lemma}

\begin{proof}
The first statement follows because
$\text{Supp}(M) = \bigcup V(\mathfrak p_i)$, see
Lemma \ref{lemma-filter-primes-in-support}.
Let $\mathfrak p \in \text{Supp}(M)$ be minimal.
The support of $M_{\mathfrak p}$ is the set
consisting of the maximal ideal $\mathfrak p R_{\mathfrak p}$.
Hence by Lemma \ref{lemma-support-point} the length
of $M_{\mathfrak p}$ is finite and $> 0$. Next we
note that $M_{\mathfrak p}$ has a filtration with subquotients
$
(R/\mathfrak p_i)_{\mathfrak p}
=
R_{\mathfrak p}/{\mathfrak p_i}R_{\mathfrak p}
$.
These are zero if $\mathfrak p_i \not \subset \mathfrak p$
and equal to $\kappa(\mathfrak p)$ if $\mathfrak p_i \subset
\mathfrak p$ because by minimality of $\mathfrak p$
we have $\mathfrak p_i = \mathfrak p$ in this case.
The result follows since $\kappa(\mathfrak p)$ has length $1$.
\end{proof}

\begin{lemma}
\label{lemma-support-dimension-d}
Let $R$ be a Noetherian local ring.
Let $M$ be a finite $R$-module.
Then $d(M) = \dim(\text{Supp}(M))$ where $d(M)$ is as in
Definition \ref{definition-d}.
\end{lemma}

\begin{proof}
Let $M_i, \mathfrak p_i$ be as in Lemma \ref{lemma-filter-Noetherian-module}.
By Lemma \ref{lemma-hilbert-ses-chi} we obtain the equality
$d(M) = \max \{ d(R/\mathfrak p_i) \}$. By
Proposition \ref{proposition-dimension} we have
$d(R/\mathfrak p_i) = \dim(R/\mathfrak p_i)$.
Trivially $\dim(R/\mathfrak p_i) = \dim V(\mathfrak p_i)$.
Since all minimal primes of $\text{Supp}(M)$ occur among
the $\mathfrak p_i$ (Lemma \ref{lemma-filter-minimal-primes-in-support}) we win.
\end{proof}

\begin{lemma}
\label{lemma-ses-dimension}
Let $R$ be a Noetherian ring. Let $0 \to M' \to M \to M'' \to 0$
be a short exact sequence of finite $R$-modules. Then
$\max\{\dim(\text{Supp}(M')), \dim(\text{Supp}(M''))\} =
\dim(\text{Supp}(M))$.
\end{lemma}

\begin{proof}
If $R$ is local, this follows immediately from
Lemmas \ref{lemma-support-dimension-d} and \ref{lemma-hilbert-ses-chi}.
A more elementary argument, which works also if $R$ is not local,
is to use that $\text{Supp}(M')$, $\text{Supp}(M'')$, and
$\text{Supp}(M)$ are closed (Lemma \ref{lemma-support-closed})
and that $\text{Supp}(M) = \text{Supp}(M') \cup \text{Supp}(M'')$
(Lemma \ref{lemma-support-quotient}).
\end{proof}










